\subsection{SNARKs: proofs of general statements}

A powerful way of thinking about a signature scheme is that it is a
\textbf{proof}. Specifically, Alice's signature is a proof that ``I
{[}the person who generated the signature{]} know Alice's private key.''
Similarly, a group signature can be thought of as a succinct proof that
``I know one of Alice, Bob, or Charlie's private keys.''

In the spirit of programmable cryptography, a \emph{SNARK} generalizes
this concept as a ``proof system'' protocol that produces efficient
proofs of \textbf{arbitrary} statements of the form:

``I know \(X\) such that \(F(X,Y) = Z\), where \(F,Y,Z\) are public,''
once the statement is encoded as a system of equations. One such
statement would be ``I know \(M\) such that
\(\operatorname{sha}(M) = Y\).''

SNARKs are an active area of research, and many different SNARKs are
known. We will focus on a particular example, PLONK ({[}plonk{]}).
